\chapter{The Morse Index Theorem}
\label{chap:dc-morse-index-theorem}

Faithful transcription of do Carmo, \emph{Riemannian Geometry}. Each numbered
statement keeps do Carmo's number, kind and (name); proofs keep his explicit
cross-references ("by the Index Lemma", "\cref{cor:dc-ch11-2-6}", "\cref{rem:dc-ch9-2-10}") so the dependency graph is recoverable from the text.

\section{Introduction}

In this chapter we wish to prove the Morse Index Theorem. This theorem relates
the number of conjugate points on a geodesic segment, counted with their
multiplicities, to the index of a certain quadratic form defined in terms of the
formula for the second variation (essentially, the expression $I_a(V,V)$).

The Morse Index Theorem is a generalization of a classical theorem of Jacobi
(see Cor. 2.9) which states that a geodesic segment minimizes the arc length
relative to "neighboring" curves with the same endpoints if and only if such a
segment has no conjugate points. The proof of the Index Theorem which we present
here uses the Index Lemma of Chapter 10.

\section{The Index Theorem}

Let $\gamma:[0,a]\to M^n$ be a geodesic. Denote by $\mathcal{V}(0,a)=\mathcal{V}$
the vector space formed by vector fields $V$ along $\gamma$, which are piecewise
differentiable and vanish at the endpoints of $\gamma$, that is,
$V(0)=V(a)=0$.

\begin{definition}[the index form]
  \label{def:dc-ch11-2-1}
  \dcref{ch11:2.1}
  \uses{prop:dc-ch4-2-5}
  The \emph{index form} of $\gamma$ is the quadratic form associated to the symmetric
  bilinear form $I_a$ defined on $\mathcal{V}$ by
  
  $$
  I_a(V,W)=\int_0^a\{\langle V',W'\rangle-\langle R(\gamma',V)\gamma',W\rangle\}\,dt,\qquad(1)
  $$
  
  where $V,W\in\mathcal{V}$ (cf. Rem. 2.10 of Chap. 9).
  
  Observe that the symmetry of $I_a$ follows from part (d) of \cref{prop:dc-ch4-2-5}.
  
  In general, given a symmetric bilinear form $B$ over a vector space
  $\mathcal{V}$, we define the \emph{index} of $B$ as the maximal dimension of all
  subspaces of $\mathcal{V}$ on which the quadratic form associated to $B$ is
  negative definite. The \emph{nullity} of $B$ is defined to be the dimension of the
  subspace of $\mathcal{V}$ formed by the elements $V\in\mathcal{V}$ such that
  $B(V,W)=0$ for all $W\in\mathcal{V}$; such a subspace is called the \emph{null space}
  of $B$. We say that $B$ is \emph{degenerate} if its nullity is strictly positive.
\end{definition}

\begin{theorem}[Morse index theorem]
  \label{thm:dc-ch11-2-2}
  \dcref{ch11:2.2}
  The index of the form $I_a$ is finite and equals the number of points
  $\gamma(t)$, $0<t<a$, conjugate to $\gamma(0)$, each counted with its
  multiplicity.
  
  Before we start the proof of the theorem, we need a few preliminary
  propositions.
\end{theorem}

\begin{proposition}[null space = Jacobi fields]
  \label{prop:dc-ch11-2-3}
  \dcref{ch11:2.3}
  An element $V\in\mathcal{V}$ belongs to the null space of $I_a$ if and only if
  $V$ is a Jacobi field along $\gamma$.
\end{proposition}
\begin{proof}
  First observe that we have the following expression for $I_a$ (compare
  Rem. 2.10 of Chap. 9):
  
  $$
  I_a(V,W)=-\int_0^a\langle V''+R(\gamma',V)\gamma',W\rangle\,dt-\sum_{j=1}^{k-1}\Big\langle\frac{DV}{dt}(t_j^+)-\frac{DV}{dt}(t_j^-),W(t_j)\Big\rangle.\qquad(2)
  $$
  
  If $V$ is a Jacobi field, then by (2), $V$ is in the null space of $I_a$.
  
  Conversely, suppose that $I_a(V,W)=0$ for all $W\in\mathcal{V}$. Let
  $0=t_0<t_1<\cdots<t_{k-1}<t_k=a$ be a subdivision of $[0,a]$ such that the
  restriction $V|[t_{j-1},t_j]$ is differentiable, $j=1,\dots,k$. Let
  $f:[0,a]\to\mathbb{R}$ be a differentiable function with $f(t)>0$ for $t\neq t_j$
  and $f(t_j)=0$, $j=0,\dots,k$. Define $W$ by
  
  $$
  W(t)=f(t)(V''+R(\gamma',V)\gamma').
  $$
  
  Then
  
  $$
  0=I_a(V,W)=-\int_0^a f(t)\|V''+R(\gamma',V)\gamma'\|^2\,dt.
  $$
  
  It follows that the integrand is zero, and therefore the restriction
  $V|(t_{j-1},t_j)$ is a Jacobi field. To see what happens at each $t_j$, choose
  $T\in\mathcal{V}$ in such a way that
  
  $$
  T(t_j)=\frac{DV}{dt}(t_j^+)-\frac{DV}{dt}(t_j^-),\quad j=1,\dots,k-1.
  $$
  
  Since
  
  $$
  0=I_a(V,T)=-\sum_{j=1}^{k-1}\Big\|\frac{DV}{dt}(t_j^+)-\frac{DV}{dt}(t_j^-)\Big\|^2,
  $$
  
  we conclude that $V$ is of class $C^1$ at each $t_j$. By the uniqueness of the
  solution to an ordinary differential equation, $V$ is $C^\infty$. Therefore $V$
  is a Jacobi field. $\square$
\end{proof}

\begin{corollary}[degeneracy and conjugate endpoints]
  \label{cor:dc-ch11-2-4}
  \dcref{ch11:2.4}
  $I_a$ is degenerate if and only if the points $\gamma(0)$ and $\gamma(a)$ are
  conjugate along $\gamma$. In this case, the nullity of $I_a$ is equal to the
  multiplicity of $\gamma(a)$ as a conjugate point.
  
  For the next proposition, as well as for the proof of the Index Theorem, we need
  some preliminary considerations.
  
  Since each point of $M$ is contained in a totally normal neighborhood and
  $\gamma([0,a])$ is compact, we can choose a subdivision
  
  $$
  0=t_0<t_1<\cdots<t_{k-1}<t_k=a
  $$
  
  of $[0,a]$ such that each $\gamma|[t_{j-1},t_j]$, $j=1,\dots,k$, is contained in a
  totally normal neighborhood. Thus each $\gamma|[t_{j-1},t_j]$ is a minimizing
  geodesic and doesn't contain any conjugate points. In what follows, such a
  subdivision will be called \emph{normal} and will be fixed until mention to the
  contrary.
  
  Let $\mathcal{V}^-(0,a)=\mathcal{V}^-$ be the vector subspace of $\mathcal{V}$
  formed from the fields $V$ such that $V|(t_{i-1},t_i)$, $i=1,\dots,k$, is a
  Jacobi field; $\mathcal{V}^-$ has finite dimension. Let $\mathcal{V}^+$ be the
  subspace of $\mathcal{V}$ consisting of vector fields $W$ such that
  $W(t_1)=W(t_2)=\cdots=W(t_{k-1})=0$.
\end{corollary}

\begin{proposition}[the splitting $\mathcal{V}=\mathcal{V}^+\oplus\mathcal{V}^-$]
  \label{prop:dc-ch11-2-5}
  \dcref{ch11:2.5}
  $\mathcal{V}$ is a direct sum $\mathcal{V}=\mathcal{V}^+\oplus\mathcal{V}^-$, and
  the subspaces $\mathcal{V}^+$ and $\mathcal{V}^-$ are orthogonal with respect to
  $I_a$. In addition, $I_a$ restricted to $\mathcal{V}^+$ is positive definite.
\end{proposition}
\begin{proof}
  Given $V\in\mathcal{V}$, let $W$ be a vector field in $\mathcal{V}^-$
  given by $W(t_j)=V(t_j)$; because $\gamma|[t_{j-1},t_j]$ does not have any
  conjugate points, such a $W$ exists and is unique. Hence
  $V-W\in\mathcal{V}^+$ and, therefore $\mathcal{V}=\mathcal{V}^+\oplus\mathcal{V}^-$.
  In addition, if $X\in\mathcal{V}^-$ and $Y\in\mathcal{V}^+$, we have
  
  $$
  I_a(X,Y)=-\sum_{j=1}^{k-1}\Big\langle 0,\frac{DX}{dt}(t_j^+)-\frac{DX}{dt}(t_j^-)\Big\rangle=0,
  $$
  
  that is, $\mathcal{V}^+$ and $\mathcal{V}^-$ are orthogonal, relative to $I_a$.
  This proves the first part of Proposition 2.5.
  
  Since $\gamma|[t_{j-1},t_j]$, $j=1,\dots,k$, are minimizing geodesics, they have
  less energy than any other paths between their endpoints. Therefore, if
  $V\in\mathcal{V}^+$, then $I_a(V,V)\geq 0$.
  
  It remains to show that $I_a(V,V)>0$ if $V\in\mathcal{V}^+-\{0\}$. Suppose, to
  the contrary, that $I_a(V,V)=0$ with $V\in\mathcal{V}^+$, $V\neq 0$. We are going
  to show that this implies that $V$ belongs to the null space of $I_a$. Indeed,
  if $W\in\mathcal{V}^-$, then $I_a(V,W)=0$, from the orthogonality above. If
  $W\in\mathcal{V}^+$, consider the inequality
  
  $$
  0\leq I_a(V+cW,V+cW)=2cI_a(V,W)+c^2 I_a(W,W),
  $$
  
  valid for all real $c$. This says that there exist real numbers $A\geq 0$ and
  $B$ such that $Ac^2+2Bc\geq 0$ for all $c\in\mathbb{R}$, which is only possible
  when $B=0$, that is, $I_a(V,W)=0$. Therefore $V$ belongs to the null space of
  $I_a$. Since the null space consists of Jacobi fields and $V$ vanishes at $t_j$,
  we conclude that $V=0$, which is a contradiction and ends the proof. $\square$
\end{proof}

\begin{corollary}[index reduces to $\mathcal{V}^-$]
  \label{cor:dc-ch11-2-6}
  \dcref{ch11:2.6}
  The index of $I_a$ is equal to the index of $I_a$ restricted to
  $\mathcal{V}^-$; in particular, the index of $I_a$ is finite. The same is true
  for the nullity of $I_a$.
  
  \emph{Proof of the Index Theorem.} It is convenient to introduce the following
  notation. If $t\in[0,a]$, denote by $\gamma_t$ the restriction of $\gamma$ to the
  interval $[0,t]$; the corresponding index form will be denoted by $I_t$, and the
  index of $I_t$ will be denoted by $i(t)$. In this manner, we define a function
  $i:[0,a]\to\mathbb{N}$, whose behavior we wish to study.
  
  Recalling that the subdivision (which is supposed fixed) of $[0,a]$ by the points
  $t_j$, $j=0,\dots,k$, was chosen in a way that $\gamma|[t_{j-1},t_j]$ is a
  minimizing geodesic, we conclude that $i(t)$ is zero in a neighborhood of $0$. In
  addition, $i(t)$ is non-decreasing, that is, if $\bar t>t$, then
  $i(\bar t)\geq i(t)$. In fact, by the definition of $i(t)$, there exists a
  subspace $\mathcal{U}\subset\mathcal{V}(0,t)$ such that $I_t$ is negative
  definite on $\mathcal{U}$ and $\dim\mathcal{U}=i(t)$. Every element
  $V\in\mathcal{U}$ can be extended to an element $\bar V\in\mathcal{V}(0,\bar t)$
  by defining $\bar V=0$ on $[t,\bar t]$. It is clear that
  $I_t(V,V)=I_{\bar t}(\bar V,\bar V)$. From the definition of the index,
  $i(\bar t)\geq i(t)$, as has been claimed.
  
  To obtain other properties of $i(t)$ we proceed in the following manner. First,
  observe that, from the definition, $i(t)$ does not depend on the choice of normal
  subdivision of $[0,a]$; we can, therefore, choose such subdivisions in a way that
  $t\in(t_{j-1},t_j)$. Next, we observe that the index of $I_t$ is the index of the
  restriction of $I_t$ to the subspace $\mathcal{V}^-(0,t)$; such a restriction
  will be again denoted by $I_t$. Then, as each element of $\mathcal{V}^-(0,t)$ is
  determined by its value at the points $\gamma(t_1),\dots,\gamma(t_{j-1})$, we have
  that $\mathcal{V}^-(0,t)$ is isomorphic to a direct sum
  
  $$
  \mathcal{V}^-(0,t)=T_{\gamma(t_1)}M\oplus\cdots\oplus T_{\gamma(t_{j-1})}M=S_j.
  $$
  
  It follows, by letting $t$ vary on $(t_{j-1},t_j)$, that the spaces
  $\mathcal{V}^-(0,t)$ are isomorphic to each other and isomorphic to $S_j$. We
  can, therefore, consider the quadratic forms $I_t$ as a family of quadratic forms
  on a fixed space $S_j$. In addition, since the elements of $\mathcal{V}^-(0,t)$
  are "broken" Jacobi fields, it follows from (2) that $I_t$ depends continuously
  on $t\in(t_{j-1},t_j)$.
  
  We can now obtain more information about $i(t)$ which will be gathered together
  in the lemmas below.
\end{corollary}

\begin{lemma}[left continuity of $i$]
  \label{lem:dc-ch11-2-7}
  \dcref{ch11:2.7}
  If $\varepsilon>0$ is sufficiently small, $i(t-\varepsilon)=i(t)$.
  
  \emph{Proof of Lemma 2.7.} Since $i(t)$ is non-decreasing,
  $i(t)\geq i(t-\varepsilon)$, for all $\varepsilon$. On the other hand, if $I_t$
  is negative definite on a subspace $\bar S\subset S_j$, with $\dim(\bar S)=i(t)$,
  then, by continuity of $I_t$, there exists $\varepsilon>0$ such that
  $I_{t-\varepsilon}$ is still negative definite on $\bar S$, hence
  $i(t-\varepsilon)\geq i(t)$. Therefore $i(t)=i(t-\varepsilon)$. $\square$
  
  Now let $d$ be the nullity of $I_t$; observe that $d=0$ if $\gamma(t)$ is not
  conjugate to $\gamma(0)$.
\end{lemma}

\begin{lemma}[right jump by the nullity]
  \label{lem:dc-ch11-2-8}
  \dcref{ch11:2.8}
  \uses{thm:dc-ch11-2-2}
  If $\varepsilon>0$ is sufficiently small, $i(t+\varepsilon)=i(t)+d$.
  
  \emph{Proof of Lemma 2.8.} We show first that $i(t+\varepsilon)\leq i(t)+d$. Indeed,
  because $\dim(S_j)=n(j-1)$, $I_t$ is positive definite on a subspace of dimension
  $n(j-1)-i(t)-d$. By continuity, $I_{t+\varepsilon}$ is still positive definite on
  this subspace, for $\varepsilon>0$ sufficiently small. Therefore
  
  $$
  i(t+\varepsilon)\leq n(j-1)-\{n(j-1)-i(t)-d\}=i(t)+d.
  $$
  
  To prove the inequality in the other direction, it is necessary to use the Index
  Lemma (\cref{lem:dc-ch10-2-2} of Chap. 10). Let $V\in S_j$, with $V(t_{j-1})\neq 0$, and
  denote by $V_{t_o}$ the "broken" Jacobi field which coincides with $V(t_i)$ at
  $t_i$, $i=1,\dots,j-1$, and which vanishes at the point $t_o\in(t_{j-1},t_j)$. We
  claim that
  
  $$
  I_{t_o}(V_{t_o},V_{t_o})>I_{t_o+\varepsilon}(V_{t_o+\varepsilon},V_{t_o+\varepsilon}).
  $$
  
  In fact, if we denote by $W_{t_o}$ (see Fig. 1) the vector field defined along
  $\gamma([0,t_o+\varepsilon])$ by:
  
  $$
  W_{t_o}(t)=V_{t_o}(t),\quad t\in[0,t_o],\qquad W_{t_o}(t)=0,\quad t\in[t_o,t_o+\varepsilon],
  $$
  
  we have, from the Index Lemma,
  
  $$
  I_{t_o}(V_{t_o},V_{t_o})=I_{t_o+\varepsilon}(W_{t_o},W_{t_o})>I_{t_o+\varepsilon}(V_{t_o+\varepsilon},V_{t_o+\varepsilon}),
  $$
  
  where the last inequality is strict, since $W_{t_o}$ is not a Jacobi field. This
  proves the assertion made.
  
  Therefore, if $V\in S_j$ and $I_t(V,V)\leq 0$, then $I_{t+\varepsilon}(V,V)<0$.
  Hence, if $I_t$ is negative definite on a subspace $\bar S\subset S_j$,
  $I_{t+\varepsilon}$ will still be negative definite on the direct sum of $\bar S$
  with the null space of $I_t$. Therefore
  
  $$
  i(t+\varepsilon)\geq i(t)+d,
  $$
  
  which, together with the previous inequality implies that
  $i(t+\varepsilon)=i(t)+d$. $\square$
  
  The information which we obtained about $i(t)$ allows us to describe $i(t)$ as a
  function which is zero in a neighborhood of the origin, is continuous on the left
  and has "step" type discontinuities at the conjugate points to $\gamma(0)$, the
  step being exactly equal to the multiplicity of the conjugate point (see Fig. 2).
  But this is precisely the statement of the Index Theorem. $\square$
\end{lemma}

\begin{corollary}[Jacobi]
  \label{cor:dc-ch11-2-9}
  \dcref{ch11:2.9}
  Let $\gamma:[0,a]\to M$ be a geodesic segment on $M$ such that $\gamma(a)$ is not
  conjugate to $\gamma(0)$. Then $\gamma$ has no conjugate points on $(0,a)$ if and
  only if for all proper variations of $\gamma$ there exists a $\delta>0$ such that
  $E(0)\leq E(s)$ for $|s|<\delta$. In particular, if $\gamma$ is minimizing,
  $\gamma$ has no conjugate points on $(0,a)$.
\end{corollary}

\begin{corollary}[conjugate points are discrete]
  \label{cor:dc-ch11-2-10}
  \dcref{ch11:2.10}
  The set of conjugate points along a geodesic is a discrete set.
\end{corollary}

\begin{remark}[Remark 2.11]
  \label{rem:dc-ch11-2-11}
  \dcref{ch11:2.11}
  The index theorem can be generalized to the case in which the points $\gamma(0)$
  and $\gamma(a)$ are replaced by submanifolds. For this see W. Ambrose [Am 2]. For
  an extremely general version of the index theorem, see S. Smale [Sm].
\end{remark}

\begin{remark}[Remark 2.12]
  \label{rem:dc-ch11-2-12}
  \dcref{ch11:2.12}
  \uses{rem:dc-ch9-2-7, cor:dc-ch11-2-6}
  The fact that the index and the nullity of $I$ are finite reflect the fundamental
  fact that the "infinite dimensional manifold" $\Omega_{p,q}=\Omega$ mentioned in
  \cref{rem:dc-ch9-2-7} can be replaced, under certain conditions, by a finite
  dimensional manifold. This was the original point of view of Morse to prove his
  theorem. For some applications (cf. the Sphere Theorem, Prop. 3.1), it is
  convenient to have this construction in mind which can be summarized in the
  following way. (For more details, see Milnor [Mi], pp. 88–92.)
  
  Let $\Omega^c$ (resp. $\overset{\circ}{\Omega}{}^c$) be the subset of $\Omega$
  formed by curves in $\Omega$ of energy $\leq c$ (resp. $<c$). It is clear that the
  curves in $\Omega^c$ are contained in a compact subset $S\subset M$. Let
  $\delta>0$ be such that given two points of $S$ at a distance less than $\delta$,
  there exists a unique minimizing geodesic joining these two points. Suppose that
  the curves in $\Omega^c$ are defined on $[0,1]$ and subdivide $[0,1]$ by the
  points $t_i$, $i=0,\dots,k$, in such a way that
  $|t_i-t_{i-1}|<\frac{\delta^2}{c}$. Let $B\subset\Omega^c$ (resp.
  $\overset{\circ}{B}\subset\overset{\circ}{\Omega}{}^c$) be the set of geodesics
  broken at $t_i$, that is, $w\in B$ if $w(0)=p$, $w(1)=q$, and the restriction
  $w|(t_{i-1},t_i)=w_i$ is a geodesic joining $w(t_{i-1})$ to $w(t_i)$. Such a
  geodesic is entirely determined by the points $w(t_i)$ since
  
  $$
  L^2(w_i)=(t_i-t_{i-1})E(w_i)<\delta^2.
  $$
  
  The correspondence
  
  $$
  w\mapsto(w(t_1),\dots,w(t_{k-1}))\in M\times\cdots\times M,\quad w\in\overset{\circ}{B},
  $$
  
  is bijective and takes $\overset{\circ}{B}$ into an open subset of
  $M\times\cdots\times M$. This allows us to introduce a differentiable structure
  (of finite dimension) on $\overset{\circ}{B}$. Observe that the tangent space
  $T_w(\overset{\circ}{B})$ to a broken geodesic $w$ corresponds to the set of
  Jacobi fields along $w$, broken at $t_i$. It is possible to show that there
  exists a homotopy
  $h_s:\overset{\circ}{\Omega}{}^c\to\overset{\circ}{\Omega}{}^c$, $s\in[0,1]$ with
  $h_0=\mathrm{ident.}$ and $h_1:\overset{\circ}{\Omega}{}^c\to\overset{\circ}{B}$,
  that is, $\overset{\circ}{B}$ is a "deformation retract" of
  $\overset{\circ}{\Omega}{}^c$.
  
  Restricting the energy $E$ to a function $\bar E$ on $\overset{\circ}{B}$, one
  verifies that the geodesics $\gamma$ of $\overset{\circ}{\Omega}{}^c$ are in the
  manifold $\overset{\circ}{B}$ and are precisely the critical points of $\bar E$.
  In addition, the index and the nullity of $I$ at $\gamma$ coincide with the index
  and the nullity, respectively, of the hessian of $\bar E$ at $\gamma$, since from
  \cref{cor:dc-ch11-2-6} they coincide with the index and the nullity, respectively, of $I$
  restricted to the broken Jacobi fields along $\gamma$.
  
  In this way, for many purposes, we can substitute the set
  $\overset{\circ}{\Omega}{}^c$ by its finite dimensional approximation
  $\overset{\circ}{B}$.
\end{remark}
